\question[3]
Eine natürliche Zahl heißt \emph{vollkommen}, wenn sie gleich der Summe ihrer echten Teiler ist,
z.B. $6 = 1 + 2 + 3$. Verwende die Funktion \texttt{echteTeilerSumme} aus der vorigen Aufgabe.
\begin{parts}
\part Schreibe eine Funktion \texttt{ist\_vollkommen(n)}, die \texttt{True} zurückgibt, wenn \texttt{n}
vollkommen ist, sonst \texttt{False}.
\part Schreibe eine Funktion \texttt{vollkommene\_zahlen(bis)}, die eine Liste aller vollkommenen Zahlen
von 1 bis \texttt{bis} zurückgibt. \texttt{vollkommene\_zahlen(500)} liefert \texttt{[6, 28, 496]}.
\end{parts}

\begin{solutionbox}{5cm}
\begin{lstlisting}
def ist_vollkommen(n):
    return echteTeilerSumme(n) == n

def vollkommene_zahlen(bis):
    ergebnis = []
    for k in range(1, bis + 1):
        if ist_vollkommen(k):
            ergebnis.append(k)
    return ergebnis
\end{lstlisting}
\end{solutionbox}
